\documentclass[11pt,a4paper]{article}
\usepackage[margin=23mm]{geometry}
\usepackage{amsmath,amssymb,booktabs,graphicx,hyperref,listings,xcolor,longtable}
\hypersetup{colorlinks=true,urlcolor=blue,linkcolor=blue}
\lstset{basicstyle=\ttfamily\scriptsize,breaklines=true,frame=single,columns=fullflexible}
\title{Finite-Range GNSS Visibility Simulation and Direction Recovery with a Distributed TART Network\\\large Technical Validation Report (Ghana pilot)}
\author{TART--GNSS collaboration\\Working technical record}
\date{7 November 2026}
\begin{document}
\maketitle
\begin{abstract}
This technical note documents the reproducible preliminary validation of GPS satellite
NORAD 39166 (NAVSTAR 68) in simulated Telescope Array Radio Telescope (TART)
visibilities. Three 24-element interferometers (Ghana, Antananarivo and Rhodes)
were simulated at GNSS L1. We distinguish the confirmed properties of the generated
measurement sets from the Ghana-only localisation tests. An orbit-blind
two-angle spherical-wave visibility fit agrees with an independent orbit prediction
to a few arcseconds in a high-signal, no-gain-error simulation; this is \emph{not}
a real-world localisation-accuracy claim.
\end{abstract}
\section{Purpose and scope}
The objective is to establish a reliable forward model before applying Bayesian
inference to multi-station GNSS localisation. A physically consistent
simulation must preserve satellite orbital elements, antenna positions, time
standards, signal phase-centre conventions, and complex visibility formation.
This work is a \emph{simulation validation}, not a validated inference on
on-sky observations. The operational \texttt{tart-gnss} stages 1--7 remain
unchanged. Scripts and run instructions are in
\texttt{docs/tabsim\_gnss/README.md} and \texttt{scripts/tabsim\_gnss\_fit.py}.
\section{Data provenance}
\begin{itemize}
\item Satellite: NORAD 39166, NAVSTAR 68 (USA 242), GPS L1
$\nu=1575.42\,\mathrm{MHz}$; saved OMM epoch
\texttt{2026-10-30T07:09:07.795296} UTC.
\item Data interval: 2026-10-29 01:10:00--01:12:01 UTC;
120 integrations, each $1.024\,\mathrm{s}$, $33$ internal samples.
\item Each MS: 24 antennas, 276 unique cross-baselines,
$120\times276=33120$ rows, one frequency channel of width $1$ MHz,
one complex correlation. The three stations have independently
generated MSs with matching displayed orbital elements.
\item Input: \texttt{tab\_sim\_\{ghana,mg-tana,za-rhodes\}.yaml}
in \texttt{configs/tart\_tabsim} of the user's local simulation working tree.
The MS and \texttt{used\_orbits.json} are deliberately not committed.
\item Software: \texttt{tabsim} v0.0.13 from
\url{https://github.com/chrisfinlay/tab-sim}; Python \texttt{casacore},
\texttt{astropy}, \texttt{skyfield}, NumPy and SciPy.
The source package is named \texttt{tabsim}, not \texttt{tab-sim}.
\end{itemize}
In an earlier installed simulator revision, a null Space-Track path prevented
the TLE simulation step from running. The resulting \texttt{0SAT} MS
contained noise but no RFI model. After upgrading, SatChecker delivered
one record 1.2487 days from the requested epoch, and the three \texttt{1SAT}
datasets contained nonzero RFI signals. The archive orbit is adequate for
\emph{internal simulation consistency}; its ephemeris accuracy against a
precision GNSS orbit has not been established.
\section{Visibility physics and conventions}
For station $s$ and antennas $i,j$, with satellite geocentric position
$\mathbf r_\mathrm{sat}(t)$, antenna positions $\mathbf r_{s,i}(t)$
in the \emph{same} celestial frame, and wavelength $\lambda=c/\nu$, define
\begin{equation}
d_{s,i}(t)=\|\mathbf r_\mathrm{sat}(t)-\mathbf r_{s,i}(t)\|,\qquad
\Delta d_{s,ij}(t)=d_{s,i}(t)-d_{s,j}(t).
\end{equation}
The simulator's satellite visibility uses phase-corrected distances
\begin{equation}
d^{\rm corr}_{s,i}=d_{s,i}+w_{s,i},\qquad
\Delta d^{\rm corr}_{s,ij}=\Delta d_{s,ij}+w_{s,ij},
\end{equation}
where $w_{s,ij}=w_{s,i}-w_{s,j}$ is the difference of
antenna phase-centre $w$ coordinates in metres. In the actual
\texttt{tabsim} implementation the sampled complex phasor is
\begin{equation}
q_{s,ij}(t)=\exp\!\left[-\frac{2\pi i}{\lambda}
\left(\Delta d_{s,ij}(t)+w_{s,ij}(t)\right)\right].
\label{eq:near}
\end{equation}
The one-second visibility is an average of internal samples, not the
phasor evaluated only at an interval midpoint:
\begin{equation}
m_{s,ij}(\boldsymbol\theta,R)=
\frac{1}{N_\mathrm{int}}\sum_{k=1}^{N_\mathrm{int}}
q_{s,ij}(t_k;\boldsymbol\theta,R).
\end{equation}
A more complete simulator forward model uses antenna-dependent apparent
amplitudes $A_{s,i}(t_k)$ inside the sum:
\begin{equation}
V^{\mathrm{sim}}_{s,ij}=
\frac1{N_\mathrm{int}}\sum_k
A_{s,i}(t_k) A_{s,j}^{*}(t_k)\,
q_{s,ij}(t_k)+\epsilon_{s,ij}.
\label{eq:vis}
\end{equation}
Our initial fit approximates the baseline-dependent amplitude with a single
complex scalar $a$; this is appropriate for a diagnostic but not a full
instrumental calibration model.

For a trial local direction $\boldsymbol\theta=(Az,El)$,
\begin{equation}
\hat{\mathbf n}_{\rm ENU}=
\begin{bmatrix}
\sin Az\cos El\\ \cos Az\cos El\\ \sin El
\end{bmatrix},\quad
\mathbf r_{\rm trial}(t)=\mathbf r_{\rm centre}(t)
+R\hat{\mathbf n}_{\rm GCRS}(t).
\label{eq:direction}
\end{equation}
The ENU basis must be rotated into the same GCRS frame as the antenna
positions. The source position is held at fixed \emph{local}
$(Az,El,R)$ during one integration; this is a short-snapshot approximation,
not a dynamics model.

For a far-field direction $(l,m,n)$ relative to the MS FIELD phase centre,
$n=\sqrt{1-l^2-m^2}$, the usual Fourier kernel depends on
$ul+vm+w(n-1)$. The first dirty-image experiment used a sign inconsistent
with the simulator's satellite visibility convention; a sign-flipped
diagnostic improved orbit agreement but did not fully resolve it.
These are conventions to \emph{test}, not arbitrary fixes to observations.
\section{Inference objective}
For complex visibilities $V_b$ and a unit-amplitude model $m_b$,
with equal weights as used in this pilot,
\begin{equation}
\chi^2(\boldsymbol\theta,R,a)
=\sum_b |V_b-a\,m_b(\boldsymbol\theta,R)|^2,
\end{equation}
where $a\in\mathbb C$ is an unconstrained nuisance amplitude. It is
eliminated analytically:
\begin{equation}
\hat a(\boldsymbol\theta,R)
=\frac{\sum_b m_b^*V_b}{\sum_b|m_b|^2},\quad
J(\boldsymbol\theta,R)=
\frac{\sum_b|V_b-\hat a m_b|^2}{\sum_b|V_b|^2}.
\label{eq:profile}
\end{equation}
The amplitude-weighted phase coherence is
\begin{equation}
C=\frac{|\sum_b m_b^* V_b|}
{\sqrt{\sum_b|m_b|^2\,\sum_b|V_b|^2}}\in[0,1].
\end{equation}
It quantifies \emph{model coherence}, not the probability that
the satellite is correctly localised. Real-data analysis will require
per-baseline weights, residual assessment, source association and gain terms.

The search centre is inferred from the MS FIELD phase centre;
the saved satellite orbit is never passed into the angle optimiser.
The optimiser uses differential evolution followed by local Nelder--Mead.
The global search has not reported success in the observed runs, whereas
local searches did. Uniqueness and coverage remain open.
\section{Measured validation results}
\subsection{MS signal components}
Mean complex visibility amplitudes (Jy-like simulated MS units) were:
\begin{center}\begin{tabular}{lrr}
\toprule
Station & $|V_{\rm RFI\,model}|$ & $|V_{\rm noise}|$\\\midrule
Ghana & $8.86947\times10^6$ & 8.76995\\
Antananarivo & $9.10358\times10^6$ & 8.76995\\
Rhodes & $8.72081\times10^6$ & 8.76995\\\bottomrule
\end{tabular}\end{center}
These signals are vastly stronger than the injected noise. They were used
to debug the geometry, not to establish representative GNSS receiver sensitivity.
\subsection{Image convention and finite-range phase test}
Using the original far-field imaging kernel, a $5^\circ$ RA rephase
produced expected $(l,m)=(-0.087058,-0.000180)$ and recovered
$(-0.088000,-0.008800)$. The displayed image discrepancy was about
$0.50^\circ$ from the \emph{original phase centre}, not necessarily
from the satellite truth. Direct comparison at the unshifted centre
against the OMM direction gave:
\begin{center}\begin{tabular}{lrrr}
\toprule
Image Fourier sign&$l$&$m$&orbit separation\\\midrule
$+1$&$-0.000600$&$-0.009000$&$1.070426^\circ$\\
$-1$&$+0.000600$&$+0.009000$&$0.339008^\circ$\\
Orbit&$-0.005316$&$+0.009077$&--\\\bottomrule
\end{tabular}\end{center}
A full spherical-wave model, sampled at 33 timestamps over
$1.024$ s and compared against \texttt{RFI\_MODEL\_DATA},
gave:
\begin{center}\begin{tabular}{lrr}
\toprule
Phase prescription & Coherence & circular phase RMS\\\midrule
$\Delta d+w_{ij}$ & 1.000000 (rounded) & $0.0479^\circ$\\
$\Delta d$ only & 0.037762 & $103.4545^\circ$\\
Conjugate of corrected model & 0.826855 & $35.0277^\circ$\\\bottomrule
\end{tabular}\end{center}
The phase RMS is a \emph{visibility phase} error; it is not a sky-coordinate
uncertainty. The test used stored snapshot $w$ throughout the short interval
and simplified amplitude weights, so its tiny residual is not exactly zero.
\subsection{Orbit-blind direction fit}
At 2026-10-29T01:10:00 UTC, the saved OMM predicts
$Az=100.950025^\circ$, $El=51.728246^\circ$,
range $21498.801$ km. For a fixed nominal range of $21500$ km
the independent visibility fit returned
$Az=100.948434157^\circ$ and $El=51.728550893^\circ$.
The azimuth/elevation differences are
$-0.001590843^\circ$ and $+0.000304893^\circ$;
the small-angle great-circle residual is about $3.7$ arcsec
using $\Delta\theta^2\simeq(\cos El\,\Delta Az)^2+\Delta El^2$.
The profile residual was $5.7500\times10^{-12}$ and
$C=0.999999999997$. This is internal consistency under idealised
conditions, not demonstrated empirical precision.
\subsection{Fixed-range sensitivity}
\begin{center}\begin{tabular}{rrrr}
\toprule
$R$ (km)& fitted $Az$ (deg)& fitted $El$ (deg)&$J$\\\midrule
18000&100.948434160&51.728550879&$5.7995\times10^{-12}$\\
20000&100.948434131&51.728550875&$5.7625\times10^{-12}$\\
21500&100.948434157&51.728550893&$5.7500\times10^{-12}$\\
23000&100.948434220&51.728550884&$5.7492\times10^{-12}$\\
26000&100.948434198&51.728550899&$5.7582\times10^{-12}$\\\bottomrule
\end{tabular}\end{center}
The fitted direction is numerically stable over this assumed range interval.
The near-constant residual implies little evidence for recovering range
from a single short-baseline station; the modest differences in $J$
must not be interpreted as a reliable range estimate. Optimiser and
floating-point sensitivity should be studied before assigning meaningful
sub-arcsecond uncertainties.
\section{Limitations, statistical requirements, and next experiments}
\begin{enumerate}
\item Fix and record the exact orbit input and software version for
all station runs. Matching displayed OMM fields are reassuring but a full
record hash is preferable for reproducibility.
\item Confirm global-optimum stability with multiple seeds and search bounds.
Test local identifiability and weak range curvature explicitly.
\item Lower simulated power and include realistic gains, phase offsets,
bandpass errors and flagged antennas. Fit \texttt{DATA}, not just the
noise-free model, when assessing empirical performance.
\item Simulate independent noise realisations and recover a covariance
$\Sigma_{Az,El}$; validate calibration by frequentist coverage of the
nominal $68\%$ intervals. Do not equate a tiny $J$ with a calibrated
uncertainty.
\item Repeat the orbit-blind direction fit for Antananarivo and Rhodes.
Only after quality control, infer a global ECEF satellite position:
\begin{equation}
p(\mathbf x\mid D)\propto p(\mathbf x)
\prod_s p\!\left(\hat{\boldsymbol\theta}_s
\mid\operatorname{AzEl}_s(\mathbf x),\Sigma_s\right).
\end{equation}
For a Gaussian local likelihood, properly wrap azimuth, account for the
elevation-dependent metric, and handle station-dependent systematic errors.
A direct joint-visibility model provides an important cross-check.
\end{enumerate}
\section{Reproducibility}
See \texttt{docs/tabsim\_gnss/README.md} for exact commands and path variables.
The companion script \texttt{scripts/tabsim\_gnss\_fit.py}
implements orbit-blind fitting and MS component inspection.
A provided output JSON is a local product; do not commit absolute machine
paths, bulky MS contents, or private credentials into the repository.
\section{Source code provenance}
Physics/convention verified against the open simulator implementation:
\url{https://github.com/chrisfinlay/tab-sim},
notably \texttt{tabsim/dask/observation.py} (distance plus antenna $w$)
and \texttt{tabsim/jax/interferometry.py} (baseline phase and averaging).
The local \texttt{tart-gnss} pipeline is:
\url{https://github.com/ifornax/tart-gnss}.
No claims about independent observational orbit accuracy or measured
TART instrumental SEFD are made here.

\section{Planned controlled noise experiment: directional uncertainty}
\label{sec:uncertainty}
The next experiment conditions on the saved Ghana satellite model snapshot
$S_b=V_b^{\mathrm{RFI\_MODEL}}$ but generates new independent receiver
noise for each trial. Let $b$ index unflagged complex baselines. Define
\begin{equation}
N_b^{(r)}=\sigma(Z_{1b}^{(r)}+iZ_{2b}^{(r)}),
\quad Z_{kb}^{(r)}\overset{\mathrm{iid}}{\sim}\mathcal N(0,1),
\quad E|N|^2=2\sigma^2.
\end{equation}
We estimate $\sigma^2$ from the stored \texttt{NOISE\_DATA} column
over available unflagged samples, and set a prescribed
\emph{single-complex-visibility} signal-to-noise ratio
\begin{equation}
\rho=\frac{\sqrt{B^{-1}\sum_b|\alpha_\rho S_b|^2}}
{\sqrt2\,\sigma},\qquad
\alpha_\rho=\frac{\rho\sqrt2\,\sigma}
{\sqrt{B^{-1}\sum_b|S_b|^2}}.
\end{equation}
The simulated observation at each repetition is
\begin{equation}
V_b^{(r)}(\rho)=\alpha_\rho S_b + N_b^{(r)}.
\end{equation}
In particular, $\rho$ is \emph{not} the coherent phased-array SNR.
We use a reproducible pseudorandom generator seed and avoid reading
the saved OMM truth until the entire fit loop has finished.

For each repetition $r$, minimise the complex-amplitude-profiled
residual in Eq.~\eqref{eq:profile} over a pointing-informed bounded
grid plus local Nelder--Mead refinement. Define the small-angle
tangent-plane error, in arcseconds,
\begin{equation}
\mathbf e^{(r)}=
3600\begin{bmatrix}
\cos(El_\star)\,\operatorname{wrap}_{[-180,180)}(Az^{(r)}-Az_\star)\\
El^{(r)}-El_\star
\end{bmatrix},
\end{equation}
where the subscript $\star$ denotes an independent orbit-predicted
reference and both coordinate differences are expressed in degrees.
This tangent approximation must be revisited for very wide residuals.

For each SNR level, split independent trials into training and
held-out evaluation halves. The sample bias and covariance from
training trials are
\begin{align}
\hat{\mathbf b}
&=\frac1{N_{\rm train}}\sum_{r\in{\rm train}}\mathbf e^{(r)},\\
\hat{\Sigma}_{EN}
&=\frac1{N_{\rm train}-1}\sum_{r\in{\rm train}}
(\mathbf e^{(r)}-\hat{\mathbf b})
(\mathbf e^{(r)}-\hat{\mathbf b})^\mathsf T.
\end{align}
These are conditional, repeated-noise error statistics; they are
\emph{not by themselves} a posterior uncertainty for an arbitrary
new observation. In-sample covariance can underestimate
out-of-distribution effects from gain errors and model misspecification.

The conventional nominal joint 68\% Gaussian ellipse contains
a held-out trial when
\begin{equation}
D^2=\mathbf e^\mathsf{T}\hat{\Sigma}_{EN}^{-1}\mathbf e
\le \chi^2_{2}(0.68)\simeq2.279.
\end{equation}
Report both uncorrected and bias-corrected held-out coverage;
for the latter replace $\mathbf e$ by
$\mathbf e-\hat{\mathbf b}$. If model bias is non-negligible,
report it separately rather than subsume it into covariance.
Accuracy is summarised by the median, 68th and 95th percentiles of
$\|\mathbf e\|$, and the root mean-square radial error.
The covariance in degrees squared is obtained by
$\Sigma_{Az,El}=D^{-1}\Sigma_{EN}(D^{-1})^\mathsf T$,
$D=\operatorname{diag}(3600\cos El_\star,3600)$.

\paragraph{Run and expected products.}
See \texttt{docs/tabsim\_gnss/README.md} and
\texttt{scripts/tabsim\_ghana\_uncertainty.py}. The recommended full
design uses $\rho\in\{3,5,10,20,50,100\}$ with 200 independent
trials per value. Outputs will be a trial-level CSV, a per-SNR
summary CSV, a covariance JSON, and a plot of angular error versus
SNR. \emph{No Monte Carlo result has yet been obtained for this
new experiment}; no plot is included as if measured.

\paragraph{Known limitations.}
The experiment conditions on one fixed source, one snapshot and the
same simulator family as the fit; the nominal 21,500 km range is held
fixed. The satellite apparent direction can change during integration;
a constant local direction parameterisation may create a systematic
error floor. The source is rescaled while the Gaussian noise variance
is kept fixed, rather than rerunning the full simulator with new
emission and instrumental gain parameters. Optimisation failures,
edge hits, and non-Gaussian multimodality must be retained and
reported. Subsequent calibration must include gain perturbations,
station-specific errors and independent snapshots before
$\Sigma_{EN}$ is used as a triangulation likelihood.


\section{Post hoc correction of the orbital time reference}
\label{sec:utc-ut1}
After the initial Monte Carlo analysis, a controlled diagnostic
compared the same saved OMM trajectory and Skyfield antenna geometry
using two interpretations of the numerical observation Julian date:
an explicit UTC interpretation and the installed simulator's
\texttt{ts.ut1\_jd} interpretation. The UTC-labelled observation
timestamp was 2026-10-29T01:10:00.000.
The simulator-compatible geometry yielded
$J=1.20662429051\times10^{-12}$, compared with
$J=3.82488077454\times10^{-7}$ for the UTC interpretation.
The predicted direction shift was
\begin{equation}
(\Delta E,\Delta N)_{\mathrm{sim-UTC}}
=(-3.54816'',+1.09650'').
\end{equation}
This explains the previously observed noise-free direction-fit
discrepancy of approximately $3.713''$: that offset is primarily a
truth-reference convention mismatch rather than measured
interferometric bias. The simulator itself has not yet been corrected
or regenerated, so the distinction between \emph{simulation-consistent}
direction and physically timestamped UTC direction must be retained.

Re-scoring the same $200$ saved Monte Carlo estimates at each SNR
against simulator-consistent truth (without rerunning fits) gave:
\begin{center}
\begin{tabular}{rrrr}
\toprule
Per-visibility SNR & UTC median error (
'$) &
 Simulator median error (
'$) & Raw held-out 68\% coverage\\
\midrule
3 & 110.72 & 111.45 & 0.55\\
5 & 63.62 & 63.97 & 0.71\\
10 & 30.54 & 30.43 & 0.74\\
20 & 16.05 & 16.48 & 0.54\\
50 & 7.05 & 6.55 & 0.71\\
100 & 4.90 & 3.35 & 0.67\\
\bottomrule
\end{tabular}
\end{center}
At SNR 100, training-sample East/North bias changed from
approximately $(-3.57'',+0.88'')$ to
$(-0.03'',-0.22'')$, and nominal raw held-out ellipse coverage
improved from 0.37 to 0.67. The covariance does not change materially
under subtraction of a constant error offset; it remains a
noise-conditional training-sample covariance and is not yet a
calibrated observation posterior. Raw coverage at SNR 20 remains
only 0.54 and should be investigated. Fitting failures and
systematic instrumental effects have not been resolved.

The reanalysis script is
\texttt{scripts/tabsim\_ghana\_rescore\_truth.py}; commands are in
\texttt{docs/tabsim\_gnss/README.md}. Future validation must fix and
test time-scale conversion at the simulator source and regenerate
a new Ghana measurement set before using a physical UTC reference
as a source of absolute position truth.

\end{document}
